\documentclass[11pt,a4paper]{article}
\usepackage[spanish,es-nodecimaldot]{babel}	% Utilizar español
\usepackage[utf8]{inputenc}					% Caracteres UTF-8
\usepackage{graphicx}						% Imagenes
\usepackage[hidelinks]{hyperref}			% Poner enlaces sin marcarlos en rojo
\usepackage{fancyhdr}						% Modificar encabezados y pies de pagina
\usepackage{float}							% Insertar figuras
\usepackage[textwidth=390pt]{geometry}		% Anchura de la pagina
\usepackage[nottoc]{tocbibind}				% Referencias (no incluir num pagina indice en Indice)
\usepackage{enumitem}						% Permitir enumerate con distintos simbolos
\usepackage[T1]{fontenc}					% Usar textsc en sections
\usepackage{amsmath}						% Símbolos matemáticos

% Comando para poner el nombre de la asignatura
\newcommand{\asignatura}{Visión por Computador}
\newcommand{\autor}{Vladislav Nikolov Vasilev}
\newcommand{\titulo}{Trabajo 1}
\newcommand{\subtitulo}{Filtrado y Detección de Regiones}

% Configuracion de encabezados y pies de pagina
\pagestyle{fancy}
\lhead{\autor{}}
\rhead{\asignatura{}}
\lfoot{Grado en Ingeniería Informática}
\cfoot{}
\rfoot{\thepage}
\renewcommand{\headrulewidth}{0.4pt}		% Linea cabeza de pagina
\renewcommand{\footrulewidth}{0.4pt}		% Linea pie de pagina

\begin{document}
\pagenumbering{gobble}

% Pagina de titulo
\begin{titlepage}

\begin{minipage}{\textwidth}

\centering

\includegraphics[scale=0.5]{img/ugr.png}\\

\textsc{\Large \asignatura{}\\[0.2cm]}
\textsc{GRADO EN INGENIERÍA INFORMÁTICA}\\[1cm]

\noindent\rule[-1ex]{\textwidth}{1pt}\\[1.5ex]
\textsc{{\Huge \titulo\\[0.5ex]}}
\textsc{{\Large \subtitulo\\}}
\noindent\rule[-1ex]{\textwidth}{2pt}\\[3.5ex]

\end{minipage}

\vspace{0.5cm}

\begin{minipage}{\textwidth}

\centering

\textbf{Autor}\\ {\autor{}}\\[2.5ex]
\textbf{Rama}\\ {Computación y Sistemas Inteligentes}\\[2.5ex]
\vspace{0.3cm}

\includegraphics[scale=0.3]{img/etsiit.jpeg}

\vspace{0.7cm}
\textsc{Escuela Técnica Superior de Ingenierías Informática y de Telecomunicación}\\
\vspace{1cm}
\textsc{Curso 2019-2020}
\end{minipage}
\end{titlepage}

\pagenumbering{arabic}
\tableofcontents
\thispagestyle{empty}				% No usar estilo en la pagina de indice

\newpage

\setlength{\parskip}{1em}

\section{\textsc{Ejercicio sobre filtros básicos}}

\noindent \textbf{USANDO LAS FUNCIONES DE OPENCV}: escribir funciones que implementen los siguientes puntos:

\begin{enumerate}[label=\textbf{\Alph*)}]
	\item \textbf{El cálculo de la convolución de una imagen con una máscara 2D. Usar una Gaussiana 2D (GaussianBlur)
	y máscaras 1D dadas por getDerivKernels). Mostrar ejemplos con distintos tamaños de máscara, valores de
	sigma y condiciones de contorno. Valorar los resultados.}
	\item \textbf{Usar la función Laplacian para el cálculo de la convolución 2D con una máscara normalizada de
	Laplaciana-de-Gaussiana de tamaño variable. Mostrar ejemplos de funcionamiento usando dos tipos de
	bordes y dos valores de sigma: 1 y 3.}
\end{enumerate}

\subsection{Apartado A}

\subsection{Apartado B}

\newpage

\section{\textsc{Ejercicio sobre pirámides y detección de regiones}}

\noindent \textbf{IMPLEMENTAR} funciones para las siguiente tareas:

\begin{enumerate}[label=\textbf{\Alph*)}]
	\item \textbf{Una función que genere una representación en pirámide Gaussiana de 4 niveles de una
	imagen. Mostrar ejemplos de funcionamiento usando bordes y justificar la elección de los parámetros.}
	\item \textbf{Una función que genere una representación en pirámide Laplaciana de 4 niveles de una imagen.
	Mostrar ejemplos de funcionamiento usando bordes.}
	\item \textbf{Construir un espacio de escalas Laplaciano para implementar la búsqueda de regiones usando el siguiente
	algoritmo:}
	\begin{enumerate}[label=\textbf{\alph*.}]
		\item Fijar sigma
		\item Repetir para N escalas
		\item Realizar supresión de no-máximos en cada escala
		\begin{enumerate}[label=\roman*.]
			\item Filtrar la imagen con la Laplaciana-Gaussiana normalizada en escala
			\item Guardar el cuadrado de la respuesta para el actual nivel del espacio de escalas
			\item Incrementar el valor de sigma por un coeficiente k.( 1.2-1.4)
		\end{enumerate}
		\item Mostrar las regiones encontradas en sus correspondientes escalas. Dibujar círculos con radio proporcional a
		la escala.
	\end{enumerate}
\end{enumerate}

\subsection{Apartado A}

\subsection{Apartado B}

\subsection{Apartado C}

\newpage

\section{Ejercicio 3}

\newpage

\section{Bonus}

\newpage

\begin{thebibliography}{5}

\bibitem{nombre-referencia}
Texto referencia
\\\url{https://url.referencia.com}

\end{thebibliography}

\end{document}

